\chapter{A imagem que fala}
\noindent\textit{Vinheta.} Um culto on-line de alta produção passa a mostrar, em realidade aumentada, uma ``presença'' que responde orações em tempo real, com voz sintética e holograma. As pessoas choram, doam, compartilham. Mas aquilo não é Deus. É \textbf{engenharia de adoração}.

\section{Mapa bíblico}
\begin{itemize}[leftmargin=*,noitemsep]
  \item \textbf{A imagem que recebe fôlego e fala} (Ap~13:15): um artefato mediando culto e coerção.
  \item \textbf{Ídolos que têm boca e não falam} (Sl~115; Is~44): a ironia antiga reaparece com esteroides digitais.
  \item \textbf{Verdadeiro culto} (Jo~4:23--24; Rm~12:1): em espírito e em verdade; corpo como sacrifício vivo.
\end{itemize}

\section{IA generativa: por que é diferente desta vez}
\begin{itemize}[leftmargin=*,noitemsep]
  \item \textbf{Produção barata de plausibilidade}: texto, voz, imagem e vídeo sob demanda.
  \item \textbf{Personalização total}: cada fiel recebe uma ``liturgia'' sob medida.
  \item \textbf{Velocidade de milagre}: respostas instantâneas criam sensação de ``poder''.
  \item \textbf{Meio que vira mensagem}: a forma tecnológica molda a expectativa espiritual.
\end{itemize}

\section{O risco da \emph{idolatria da conveniência}}
Quando o culto vira produto e a presença de Deus é confundida com uma \textbf{experiência fluida}, a conveniência cobra \textbf{lealdade}. É mais fácil obedecer à interface do que ao Senhor. \textbf{Lei 1 do Engano Digital}: \emph{conveniência sempre cobra lealdade}.

\begin{nicbox}
\textbf{NIC aplicado}\\
\textbf{Narrativa}: ``Deus quer que você se sinta sempre bem'' \emph{vs.} ``Deus quer torná-lo santo''.\\
\textbf{Identidade}: consumidor de conteúdos religiosos \emph{vs.} adorador que se oferece a Deus.\\
\textbf{Controle}: plataformas como novos ``átrios'' que abrem e fecham portas de visibilidade e monetização.
\end{nicbox}

\section{Mini-exegese: o que \emph{não} forçar no texto}
Ap~13 não exige um \emph{modelo tecnológico} específico. Exige discernir \textbf{adoração + coerção}. IA, hologramas e realidades mistas podem ser \emph{meios} de idolatria, mas a raiz é sempre a \textbf{lealdade do coração}. 

\section{Ferramenta — Auditoria de culto (30 minutos para líderes)}
\begin{checklist}[title={Perguntas para o conselho/produção}]
\begin{itemize}[leftmargin=*,noitemsep]
  \item Nosso culto facilita \textbf{participação} ou apenas \textbf{consumo}? 
  \item Usamos tecnologia para \textbf{clareza e serviço} ou para \textbf{deslumbrar e prender}? 
  \item Nossa liturgia forma \textbf{adoradores} (Jo~4) ou \textbf{clientes}?
  \item O que permanece se perdermos \textbf{luzes, telas e automação} por 30 dias?
\end{itemize}
\end{checklist}

\section{Práticas para a igreja local}
\begin{checklist}
\begin{itemize}[leftmargin=*,noitemsep]
  \item \textbf{Elementos irredutíveis}: leitura pública da Palavra, oração congregacional, ceia/batismo, cânticos bíblicos, generosidade.
  \item \textbf{Tecnologia serva}: projetor e som para \emph{clareza}; sem pirotecnia vazia.
  \item \textbf{Liturgias caseiras}: culto doméstico semanal sem tela (Salmo, breve Palavra, oração).
  \item \textbf{Transparência}: publique orçamento e critérios de produção; corte o que não serve à edificação.
\end{itemize}
\end{checklist}

\section{Perguntas para grupos pequenos}
\begin{itemize}[leftmargin=*,noitemsep]
  \item Em que momentos você confunde \textit{experiência} com \textit{presença de Deus}?
  \item Que tecnologia precisa ser recolocada no seu culto como \textbf{serva}, não como \textbf{senhora}?
  \item O que sua igreja manteria se perdesse a infraestrutura por um mês?
\end{itemize}

\bigskip
\noindent\textit{Próximo capítulo: Profetas do algoritmo --- recomendação, bolhas e fabricação de consenso.}
